% Online Appendix to Section 7.5 (sec:coalition). Numbers: paper/results/coalition.json.
\section{Coalition rule: data and further results}\label{oa:coalition}

\paragraph{Supply and reserves measures of the supplier route.} An economy passes the supplier test
if, in at least one Upstream or Midstream stage of a headline technology, it holds at least 5
percent of the world total of one of the following supply measures, or at least 10 percent of world
reserves of a resource the stage extracts. Supply: exports of the stage's HS6 basket
\citep[2023;][]{growthlab2019trade}; oil, gas and coal production \citep[2025;][]{ei2026statistical};
uranium mine production (2024, \citealp{oecdnea2026uranium}); mine production and refining of the
minerals a technology uses \citep[2025;][with manganese, zinc and platinum-group metals from the
Energy Institute]{iea2026cmdata}; and clean-technology manufacturing capacity (2024). Reserves: the
critical-mineral reserves of \citet{usgs2026mcs} at the end of 2025 (cobalt, lithium, graphite, rare
earths, copper, manganese, nickel, zinc and platinum-group metals; the report's erratum puts Brazil's
rare earths at 11 Mt), proved oil, gas and coal reserves at the end of 2020 \citep{ei2026statistical}
and identified uranium resources recoverable below USD 260/kgU \citep{oecdnea2026uranium}. Holders
whose reserves USGS does not quantify (entries NA, W or s) count as zero, so world shares are
measured on quantified holdings. The route reads the same reserves tables as the Reserves category
of the Energy Security pillar (Section~\ref{sec:es}). The 10 percent bar binds for a small group:
six economies that are not China-leaning hold between 5 and 10 percent of a world reserve without
otherwise qualifying (Madagascar 8.7 percent of graphite, Mozambique 8.1 percent, Kuwait 5.9 percent
of oil, Tanzania 5.8 percent, Niger 5.7 percent of uranium, New Caledonia 5.1 percent of nickel), as
do Iraq, Turkmenistan and the United Arab Emirates, which are China-leaning.

\paragraph{GDP floor.} The floor of USD 100 billion (current dollars, latest year, 2025 for every
member) excludes 18 U.S.-leaning economies, among them six EU members and small European and Pacific
states, and 16 members of the five agreements, the largest Lithuania (USD 95 billion), which enters
through the EU node, and Panama (USD 91 billion). Puerto Rico inherits the U.S. position and exceeds
the floor but has no vote of its own and is not a member.

\paragraph{Supply cover of the coalition.} Across the 18 Upstream and Midstream stages with an HS6
basket, the 63 economies of the coalition (the 57 members and the six EU members below the GDP floor) account for a median 70.7 percent of world exports, from
45.2 percent in Wind Upstream to 88.4 percent in Batteries Upstream, and China for a median 17.6
percent, with a maximum of 40.9 percent in Batteries Midstream. The coalition's share exceeds China's
in all 18 stages, but in two Upstream stages the excluded China-leaning suppliers other than China
hold large shares: 31.3 percent of Wind Upstream exports (coalition 45.2) and 47.6 percent of Nuclear
Upstream exports (coalition 50.0). Coalition economies mine 81.9 percent of the world's cobalt, 79.4
percent of its nickel, 60.0 percent of its lithium and 54.5 percent of its copper, but 3.6 percent of
its graphite and 14.3 percent of its rare earths; they refine 17.4 percent of cobalt, 23.2 percent of
copper, 24.6 percent of lithium, 5.5 percent of graphite and 12.6 percent of rare earths, against
China's 75.6, 46.8, 73.7, 94.4 and 83.9 percent; nickel refining is 55.7 percent in the coalition.
Coalition economies hold 75.8 percent of quantified world cobalt reserves, 77.3 percent of lithium,
85.5 percent of manganese, 78.8 percent of nickel and 83.9 percent of platinum-group metals, but 33.3
percent of rare earths (China 58.4) and 35.7 percent of graphite (China 32.3). In oil and gas the
excluded China-leaning suppliers other than China hold about as much as the coalition or more (35.3
against 35.7 percent of oil reserves, 43.4 against 30.8 percent of gas).

\paragraph{Dyadic alignment among members.} Counting EU members individually, the 1,596 pairs of
members split 29.3, 64.8 and 5.8 percent across $A_{ij}=1$, $\tfrac12$ and 0. The three China-leaning
partners have $A_{ij}=0$ with every core member and $A_{ij}=1$ with China and Russia.

\paragraph{Sensitivity of the blocs.} With session 78 alone (192 economies; Spearman 0.966 with the
headline, 89.1 percent in the same bloc), New Zealand falls just
below the U.S.-leaning cut ($s=0.583$ against 0.584) and becomes a partner through IPEF, Saudi Arabia
becomes China-leaning and loses its supplier route, Myanmar becomes nonaligned and enters as a
supplier, and India becomes China-leaning but remains a partner through IPEF and the MSP. With the
final-passage series (sessions 76--80, 193 economies, 89.6 percent in the same bloc) the cuts move to
0.446 and 0.055: China-leaning economies cluster more tightly around China in final-passage votes,
and Russia ($s=0.117$), Kazakhstan (0.056) and Myanmar (0.163) would enter the coalition as major
suppliers; Malaysia becomes China-leaning but stays through IPEF.

\paragraph{The earlier hand-picked list.} Of the 23 members fixed by hand in an earlier version, the
United States, Canada, Japan, Australia, Korea, the United Kingdom, Germany, France, Italy, Norway and
Denmark are core members, and India, Mexico, Brazil, Saudi Arabia, South Africa, Indonesia, Viet Nam,
Argentina, Morocco and Chile are partners. The United Arab Emirates (largest supply share 4.2
percent, oil reserves 5.6 percent) and Kenya leave. The 16 new partners are Colombia, Costa Rica, the
Democratic Republic of the Congo, the Dominican Republic, Gabon, Ghana, Guatemala, Malaysia, Namibia,
Oman, Peru, the Philippines, Qatar, Singapore, Thailand and T\"urkiye.
